\chapter{Heat Stroke}

\section*{Definition}
Heat stroke is a life-threatening condition characterized by core body temperature above 40 C (104 F) and central nervous system dysfunction. It represents the most severe form of heat-related illness and requires emergency medical treatment.

\section*{Pathophysiology}
Heat stroke occurs when thermoregulatory mechanisms fail under extreme heat stress. The body's core temperature rises uncontrollably, leading to denaturation of proteins, cellular damage, and a systemic inflammatory response. This triggers multi-organ dysfunction including encephalopathy, acute kidney injury, liver failure, disseminated intravascular coagulation, and cardiovascular collapse.

\section*{Causes}
\begin{itemize}
\item Exertional: strenuous physical activity in hot, humid conditions (common in athletes, military recruits, laborers)
\item Non-exertional (classic): passive exposure to extreme heat in vulnerable populations (elderly, infants, chronic medical conditions)
\item High ambient temperature and humidity
\item Inadequate hydration
\item Poor ventilation or lack of air conditioning
\item Obesity
\item Cardiovascular disease
\item Medications: diuretics, anticholinergics, beta-blockers, antihistamines
\item Alcohol use
\item Drug use (ecstasy/MDMA, amphetamines, cocaine)
\end{itemize}

\section*{When to See a Doctor}
\begin{itemize}
\item Core body temperature above 40 C (104 F)
\item Altered mental status: confusion, agitation, delirium, seizures, coma
\item Hot, dry, red skin (non-exertional type) or profuse sweating (exertional type)
\item Rapid, strong pulse
\item Headache, dizziness, lightheadedness
\item Nausea and vomiting
\item Rapid, shallow breathing
\item Hypotension or shock
\item Multi-organ failure
\end{itemize}

\section*{Related Symptoms}
\begin{itemize}
\item Heat exhaustion (precursor, milder form)
\item Heat cramps
\item Heat syncope (fainting from heat)
\item Dehydration
\item Rhabdomyolysis (muscle breakdown)
\item Acute kidney injury
\end{itemize}

\section*{Self-Care/Home Management}
\begin{itemize}
\item This is a medical emergency - CALL EMERGENCY SERVICES
\item Move the person to a cool, shaded area
\item Remove excess clothing
\item Cool the body aggressively: immerse in cool water, apply ice packs to neck, armpits, and groin
\item Fan the person while misting with cool water
\item Do NOT give fluids if the person is confused or unconscious
\item Monitor breathing and consciousness
\end{itemize}

\section*{How Your Doctor Will Evaluate You}
\begin{itemize}
\item Core body temperature measurement (rectal thermometer)
\item Clinical diagnosis based on hyperthermia and CNS dysfunction
\item Complete blood count, electrolytes, renal and liver function
\item Coagulation profile (disseminated intravascular coagulation)
\item Arterial blood gas analysis
\item Creatine kinase levels (rhabdomyolysis)
\item Cardiac monitoring for arrhythmias
\end{itemize}

\section*{Treatment Options}
\begin{itemize}
\item Rapid cooling to target temperature of 39 C (102 F) within 30 minutes
\item Cold water immersion (most effective method)
\item Evaporative cooling with mist and fans
\item Ice packs to neck, axillae, groin
\item Cold IV fluids (4 C normal saline)
\item Benzodiazepines for shivering or seizures
\item Supportive care for multi-organ dysfunction
\item ICU admission for severe cases
\end{itemize}
